\section{Pattern 2：Generator（生成器）}

Tool Wrapper 负责\textbf{应用知识}，而 Generator 则负责\textbf{强制输出一致性}。如果你发现 Agent 每次运行时生成的文档结构都不同，Generator 通过编排一个填空式流程来解决这个问题。

\begin{importantbox}{Generator 的关键目录结构}
Generator 利用两个可选目录：
\begin{itemize}
    \item \texttt{assets/}：存放输出模板
    \item \texttt{references/}：存放风格指南
\end{itemize}
\texttt{SKILL.md} 中的指令充当项目经理角色，告诉 Agent 加载模板、读取风格指南、向用户询问缺失变量，然后填充文档。
\end{importantbox}

\subsection{工作原理}

Generator 的执行流程是严格有序的：

\begin{enumerate}
    \item 加载 \texttt{assets/} 中的输出模板
    \item 读取 \texttt{references/} 中的风格指南
    \item 向用户询问缺失的变量
    \item 按照模板和风格指南填充文档
\end{enumerate}

关键在于：\texttt{SKILL.md} 文件本身不包含实际的布局或语法规则，它只负责\textbf{协调资产的检索}并强制 Agent 按步骤执行。

\subsection{适用场景}

\begin{itemize}
    \item 生成可预测的 API 文档
    \item 标准化 commit message
    \item 脚手架式的项目架构搭建
\end{itemize}

\begin{warningbox}{常见误区}
不要把模板内容直接写在 \texttt{SKILL.md} 中。Generator 模式的精髓在于\textbf{关注点分离}：指令文件负责流程编排，模板和风格指南作为独立资产存储在各自的目录中。这样更换模板或风格指南时不需要修改 skill 逻辑。
\end{warningbox}

\subsection{本章小结}

Generator 通过模板 + 风格指南 + 填空流程的组合，确保 Agent 每次生成的文档结构一致。指令与资产的分离使得 skill 具有高度可复用性。

